% Source-pinned comparison supporting the current Foundation endpoint adapter.
\section{The exact current Foundation endpoint and its historical interface}
\label{sec:foundation-endpoint-comparison}

This section distinguishes the current Foundation theory from the
historical vendored theory used by the PA196 calibration files.  The
distinction is necessary: their finite arithmetic axiom sets are not
literally equal, and their named $\mathsf I\Sigma _1$ bases use different
syntactic induction classes.  Their arithmetic syntax and ten-rule
proof encodings have the same numerical constructor equations.

\paragraph{Source pins and scope of the receipt.}
The ProveIt snapshot is commit
\texttt{1c9524d8bf2e0a27b9a16ff622359927f536e663}.
Its root \texttt{lakefile.toml} uses the Foundation source directory
\texttt{lib/FormalizedFormalLogic-Foundation}; the gitlink pins that
submodule to Foundation commit
\texttt{32e1a0956a8622fad067328ca1959729a7634428}.
The historical sources are the files in
\texttt{Computability/BusyBeaver/Research/PA/jamesyc-upstream/pa196/}
\texttt{vendor/Foundation/} at the ProveIt commit above.
The source manifests shipped beside this note record each fetched
blob identifier.  The historical Foundation namespace is
\texttt{FFL}; the current namespace is \texttt{LO}.  Most historical
\texttt{FirstOrder/Arithmetic/Bootstrapping} modules became current
\texttt{FirstOrder/Bootstrapping} modules.

The executable \texttt{compare\_endpoint.py} first checks all 69 saved
files against their pinned Git blob hashes and then compares 119 displayed
constructor equations, recursive-equation statements, recognizer
clauses, and finite axiom templates.  Its
\texttt{comparison\_receipt.json} retains both source strings for
each check.  All 119 checks agree after the explicitly recorded
normalizations: whitespace, namespace rename, quantifier-notation
rename, omission of a redundant local implicit binder in the
\texttt{vecToNat} declaration, and removal of the true conjunct
in \texttt{InductionR(True)}.  The changed arithmetic leaf is
recorded separately.  This is a reproducible source comparison,
not a new Lean build or an arithmetic proof checker.  The mathematical
identifications justified by these equations are proved below.

\subsection{Identical arithmetic syntax and proof encodings}

Both sources use the square pairing
\[
 P(a,b)=\begin{cases}b^2+a,&a<b,\\a^2+a+b,&b\le a.\end{cases}
\]
This is the source's \texttt{Nat.pair}, not the Cantor pairing used
by the numerical machine calculus.  Tuples are right associated:
$\langle a,b,c\rangle=P(a,P(b,c))$.
Vectors use $\mathrm{nil}=0$ and $a::v=P(a,v)+1$;
\texttt{Matrix.vecToNat} is the right fold of this same constructor.
Finite sequents use binary set coding: membership is a bit test,
and inserting $i$ leaves the code unchanged if that bit is set,
otherwise adds $2^i$.  The displayed definitions of
\texttt{Bit}, \texttt{bitInsert}, and \texttt{bitRemove} agree.
These statements are found respectively in
\texttt{Arithmetic/IOpen/Basic}, \texttt{Arithmetic/HFS/Vec},
\texttt{Vorspiel/Matrix}, and \texttt{Arithmetic/Exponential/Bit}.

The arithmetic function and relation alphabets have the following
identical arity/code pairs, as explicitly checked in both versions
of \texttt{Bootstrapping/Syntax/Language.lean}:
\[
 \begin{array}{c|rrrrrr}
 \text{symbol}&0&1&+&\times&=&<\\\hline
 \text{arity}&0&0&2&2&2&2\\
 \text{code}&0&1&0&1&0&1
 \end{array}
\]
The tags, arities, and argument-vector order are identical:
\[
 \begin{array}{rcl}
 \operatorname{Bvar}(i)&=&P(0,i)+1,\\
 \operatorname{Fvar}(i)&=&P(1,i)+1,\\
 \operatorname{Func}(k,f,v)&=&\langle2,k,f,v\rangle+1,\\
 \operatorname{Rel}(k,r,v)&=&\langle0,k,r,v\rangle+1,\\
 \operatorname{NRel}(k,r,v)&=&\langle1,k,r,v\rangle+1,\\
 \top&=&P(2,0)+1=7,\qquad \bot=P(3,0)+1=13,\\
 A\land B&=&\langle4,A,B\rangle+1,\\
 A\lor B&=&\langle5,A,B\rangle+1,\\
 \forall A&=&P(6,A)+1,\qquad\exists A=P(7,A)+1.
 \end{array}
\]
Negation exchanges truth/false, positive/negative atoms,
conjunction/disjunction, and universal/existential recursively;
source implication is $\neg A\lor B$.  The current notation
$\forall^0,\exists^0$ and the historical notation
$\forall^1,\exists^1$ designate the same constructors
\texttt{Semiformula.all}, \texttt{Semiformula.exs}; the superscript
change does not shift a de Bruijn index.

The receipt also compares every constructor equation for term
substitution, free-variable shift, bound-variable shift, formula
negation/shift/substitution, term bound-variable height, and formula
bound-variable height.  In particular, $\operatorname{qVec}$
prepends $\operatorname{Bvar}(0)$ to the bound-shifted replacement
vector, and
\[
 \operatorname{Free}(A)=
 \operatorname{Subst}([\operatorname{Fvar}(0)],
                         \operatorname{Shift}(A)).
\]
The two versions have the same well-formed-term and well-formed-
formula clauses.  For syntax validity, compute the bounded table
of all child codes up to the input code.  For substitution and
shift, generate the actual finite dependency queries
$(r,\mathbf u)$, where $r$ is the source subterm or subformula
code and $\mathbf u$ carries the entire replacement vector,
arity, and fresh-variable parameters.  A quantifier child carries
its lifted vector $\operatorname{qVec}(w)$; a source formula's
atomic case invokes the already completed term operation.
Same-operation edges strictly decrease the source code, even when
the ancillary parameter or output codes increase.  Evaluate the
finite dependency tree from its leaves.  The bounded conjunction
of valid generated queries and equality of both returned results
at all evaluated slots is a quantifier-free invariant.  It never
quantifies over all replacement vectors inside the induction
hypothesis.  Vector maps use the analogous finite list scan with
the vector held in the query packet; finite-set maps use a bounded
bit scan.  This gives primitive recursive trace constructors and
weak-base correctness proofs with the initial parameters free.
Thus the operations agree on all valid internal inputs in every
weak-base model, not only on standard test examples.

All proof codes are a right-associated tuple plus one.  The first
field is the conclusion sequent $s$ and the next field is a tag:
\begin{center}\small
\begin{tabular}{c|l|p{0.59\linewidth}}
Tag & Remaining fields & Local rule condition, in addition to valid child proofs\\\hline
0 & $p$ & $p,\neg p\in s$.\\
1 & $0$ & $\top\in s$.\\
2 & $p,q,d_p,d_q$ & $p\land q\in s$; child sequents are $\{p\}\cup s$ and $\{q\}\cup s$.\\
3 & $p,q,d$ & $p\lor q\in s$; child sequent is $\{p,q\}\cup s$.\\
4 & $p,d$ & $\forall p\in s$; child sequent is $\{\operatorname{Free}(p)\}\cup\operatorname{setShift}(s)$.\\
5 & $p,t,d$ & $\exists p\in s$, $t$ is a term; child sequent is $\{\operatorname{Subst}([t],p)\}\cup s$.\\
6 & $d$ & Child sequent is a subset of $s$.\\
7 & $d$ & $s$ is the free-variable shift of the child sequent.\\
8 & $p,d_1,d_2$ & Child sequents are $\{p\}\cup s$ and $\{\neg p\}\cup s$.\\
9 & $p$ & $p\in s$ and $p$ belongs to the chosen theory's arithmetic axiom predicate.
\end{tabular}
\end{center}
Every row additionally requires that every member of $s$ is a
well-formed formula.  The displayed \texttt{Derivation.Phi}
definition, including all ten side conditions, is identical after
whitespace normalization.  The receipt also compares
$y\in\operatorname{setShift}(s)\leftrightarrow
\exists x\in s\,(y=\operatorname{Shift}(x))$.
Finally, both predicates are
\[
 \operatorname{Proof}_T(d,A)\quad\equiv\quad
 \operatorname{Derivation}_T(d)\land
 \operatorname{fstIdx}(d)=\{A\}.
\]
Consequently the only theory-dependent constructor is tag9.

\subsection{The unrestricted induction recognizer agrees exactly}

The current \texttt{InductionSchemeDelta1.lean} defines
\[
 \begin{split}
 \operatorname{InductionUnivR}(p)\ \equiv\ 
 \exists m,b\le p\;[\;&p=\operatorname{Alls}(b,m)
 \land\operatorname{IsUFormula}(b)
 \land\operatorname{Shift}(b)=b
 \land\operatorname{bv}(b)=m\\
 &\land\exists K\le\operatorname{Subst}(\operatorname{fvarVec}(m),b)
   \;(\operatorname{IsSemiformula}_1(K)\land
       \operatorname{Subst}(\operatorname{fvarVec}(m),b)
          =\operatorname{indBodyVal}(K))].
 \end{split}
\]
The historical \texttt{Incompleteness/Definability.lean} defines
\texttt{InductionR(S)} by the same formula with the additional
conjunct $S(K)$.  Taking $S(K)$ to be true removes exactly that
conjunct.  No witness bound is enlarged or discarded: in particular
the bound on $K$ is the computed opened-body code, not $p$.

The auxiliary recurrences agree literally:
\[
 \operatorname{Alls}(b,0)=b,\quad
 \operatorname{Alls}(b,m+1)=\forall\operatorname{Alls}(b,m),
\]
\[
 \operatorname{fvarVec}(0)=0,\qquad
 \operatorname{fvarVec}(m+1)=
       \operatorname{concat}(\operatorname{fvarVec}(m),
                              \operatorname{Fvar}(m)).
\]
Both induction-body functions are, with the literal right association,
\[
 \operatorname{indBodyVal}(K)=
 K[0]\longrightarrow
   \bigl((\forall(K\longrightarrow K[\#0+1]))
                   \longrightarrow\forall K\bigr).
\]
There is a presentational difference in the two definitions of the
singleton substitution vectors.  The current source writes the
typed \texttt{SemitermVec.val} of the quoted zero and successor
terms; the historical source writes the cast of
\texttt{substCode} of those same fixed terms.  Its
\texttt{substCode\_eq} gives the equality, and the common
\texttt{vecToNat}/term-constructor equations identify the same
fixed numeral in every arithmetic model.  Thus the presentational
difference does not alter the function.  The witness triple
$(m,b,K)$ is preserved by the current-to-historical interface.

Both source theories define PA as their respective finite
\texttt{PeanoMinus} plus all universally closed
\texttt{succInd} formulas.  The current arithmetic PA axiom
predicate is assembled by
\texttt{Theory.\(\Delta_1\).add PeanoMinus.delta1
InductionScheme.delta1\_univ}.  The historical presentation uses
the same union with \texttt{InductionR(True)}.  The preceding
equations identify the induction part in every weak-base model;
the finite part still needs the following adapter.

\subsection{The one changed finite axiom and the base-theory inclusion}

The equality component consists of the same nine templates:
reflexivity, symmetry, transitivity; function congruence for
$0,1,+,\times$; and relation congruence for $=,<$.  The source
definitions of all five template families are identical after
the quantifier-notation rename.  Sixteen of the seventeen named
\texttt{PeanoMinus} arithmetic formulas are identical.  The sole
changed template is
\begin{align*}
 A_{\rm old}&=\forall x\forall y\,
   (x<y\longrightarrow\exists z\,(z<y+1\land x+z=y)),\\
 A_{\rm cur}&=\forall x\forall y\,
   (x<y\longrightarrow\exists z\,(x+z=y)).
\end{align*}
The historical notation $\exists z<^+y$ expands to
$\exists z\,(z<y+1\land\cdots)$; it does not mean $z<y$.
Existential conjunction projection proves $A_{\rm old}\to A_{\rm cur}$
by a fixed first-order derivation.  Thus a current PA proof can be
compiled to the historical arithmetic presentation by replacing
this one tag9 leaf by that fixed derivation; all other finite and
induction leaves retain their formula codes.  Rebuilding each of
the ten proof constructors with translated children is primitive
recursive.  The bounded conjunction of the correctness statements
for all processed child codes proves its correctness by
quantifier-free induction in a definitional primitive-recursive
extension of the weak base.

For reuse of the historical all-model syntax lemmas, one also needs
the reverse finite-base implication.  Current PA-minus proves it
without any induction.  Given $x<y$ and $x+z=y$, its $0\le x$
axiom means $0=x$ or $0<x$.  In the first case, commutativity and
$x+0=x$ give $z=y$.  In the second case, order preservation under
addition gives $z<y$.  Independently $0<1$, addition order
preservation, commutativity, and $x+0=x$ give $y<y+1$.
In either case $z<y+1$, using transitivity when necessary.
Pair this inequality with $x+z=y$, introduce the existential, and
close $x,y$.  This is a fixed finite derivation of $A_{\rm old}$
from current PA-minus.  Equality congruence supplies the indicated
rewrites.  Hence current PA-minus proves every historical
PA-minus axiom, and conversely.

The names $\mathsf I\Sigma _1$ nevertheless do not have literally
the same definitions.  Historical \texttt{Schemata.lean} uses
$\mathcal B[<].\operatorname{PrenexHierarchy}(\Sigma,1)$;
the current file uses
$\operatorname{Arithmetic.Hierarchy}(\Sigma,1)$.
The required inclusion is explicit.  Historical
\texttt{BoundingHierarchy.PrenexHierarchy.sigma\_succ\_iff}
and \texttt{zero\_iff\_bounded} say that an old prenex
$\Sigma_1$ formula is one existential quantifier over an old
$\mathcal B[<].\operatorname{Closure}$ matrix.  That closure is
generated by truth, falsehood, positive/negative atoms,
conjunction, disjunction, and bounded universal/existential
quantification with a \texttt{Positive} bound term.
Since $\mathcal B[<]$ is the singleton less-than operator, each of
these eight clauses is exactly a current
$\operatorname{Hierarchy}(\Pi,0)$ constructor.  Structural
induction maps the old matrix witness to a current level-zero
witness, and the current \texttt{Hierarchy.sigma} constructor
adds its outer existential.  The formula itself is unchanged.
Therefore every old prenex-$\Sigma_1$ induction axiom belongs to
the current induction class.  Together with the fixed finite-base
derivation above, every current $\mathsf I\Sigma _1$ model is a
model of the historical $\mathsf I\Sigma _1$ presentation.

This proves the precise interface needed by the calibration:
the current source predicate is handled by the literal ten-rule
translator with one finite arithmetic leaf replacement, and the
historical syntax lemmas apply over the current weak base by the
stated inclusion.  It does not assert equality of the old and new
PA proof predicates or of their defining formula texts.  The
companion Foundation-to-ordinary-PAeq proof composes this interface
with its explicit language/proof compiler to obtain the endpoint
\[
 \operatorname{Proof}_{\mathrm{PA}_{\rm current}}(d,13)
 \ \Longrightarrow\
 \operatorname{Proof}_{\mathrm{ordinaryPAeq}}(D_c(d),13).
\]
